\documentclass[conference]{IEEEtran}
\IEEEoverridecommandlockouts
\usepackage{cite}
\usepackage{amsmath,amssymb,amsfonts}
\usepackage{algorithm}
\usepackage{algorithmic}
\usepackage{graphicx}
\usepackage{textcomp}
\usepackage{xcolor}
\usepackage{booktabs}
\usepackage{multirow}
\usepackage{float}
\usepackage{caption}
\usepackage{pgfplots}
\pgfplotsset{compat=1.18}
\def\BibTeX{{\rm B\kern-.05em{\sc i\kern-.025em b}\kern-.08em
    T\kern-.1667em\lower.7ex\hbox{E}\kern-.125emX}}

\begin{document}

\title{BodhyaAI: An Explainable and Generative AI System for Cognitive-Aware Academic Mentorship}

\author{
\IEEEauthorblockN{Ms. Revathi S}
\IEEEauthorblockA{\textit{Assistant Professor, Department of AIML} \\
\textit{CMR Institute of Technology} \\
Bengaluru, India \\
revathis@cmrit.ac.in}
\and
\IEEEauthorblockN{Poorna Teja Reddy K}
\IEEEauthorblockA{\textit{Student, Department of AIML} \\
\textit{CMR Institute of Technology} \\
Bengaluru, India \\
pore22csaiml@cmrit.ac.in}
\and
\IEEEauthorblockN{Manikanta}
\IEEEauthorblockA{\textit{Student, Department of AIML} \\
\textit{CMR Institute of Technology} \\
Bengaluru, India \\
mas22csaiml@cmrit.ac.in}
\and
\IEEEauthorblockN{Sree Ranganath S}
\IEEEauthorblockA{\textit{Student, Department of AIML} \\
\textit{CMR Institute of Technology} \\
Bengaluru, India \\
srs22csaiml@cmrit.ac.in}
\and
\IEEEauthorblockN{C Karthik}
\IEEEauthorblockA{\textit{Student, Department of AIML} \\
\textit{CMR Institute of Technology} \\
Bengaluru, India \\
cka22csaiml@cmrit.ac.in}
}

\maketitle

\begin{abstract}
Traditional academic mentoring systems face critical limitations in scalability, timeliness, and objectivity, often relying on manual risk assessment and sporadic student-mentor interactions that fail to identify at-risk students proactively. Existing educational AI approaches lack interpretability, providing opaque predictions without explaining underlying causes, and rarely integrate psychological profiling with academic analytics, limiting their effectiveness in personalized intervention design.

This paper addresses these challenges through \textbf{BodhyaAI}, an Explainable and Generative AI-driven mentorship platform that unifies predictive analytics, cognitive profiling, and generative reasoning for transparent, data-driven student support. The proposed approach integrates four microservices: a \textbf{Risk Analysis Service} employing interpretable ensemble models (XGBoost, Random Forest) to predict academic and dropout risks across 21 behavioral, academic, and socioeconomic features; a \textbf{Cognitive Profiling Service} applying the Big Five Inventory (BFI-44) to assess personality dimensions (OCEAN) influencing learning outcomes; an \textbf{Explainable AI (XAI) Service} utilizing SHAP and LIME algorithms to generate feature-level explanations, ensuring transparency and mentor trust; and a \textbf{Generative MentorBot Service} powered by the Google Gemini API within a Retrieval-Augmented Generation (RAG-lite) framework, with architectural support for local model deployment (Phi-3-mini, Llama) to deliver context-aware mentoring recommendations. These services are orchestrated through a Node.js backend and React-based frontend, providing real-time dashboards, explainability visualizations, and interactive GenAI interfaces.

Trained and validated on an academic dataset of 50,000 students using five-fold stratified cross-validation, BodhyaAI achieves 95.23\% test accuracy and a weighted F1-score of 0.95 in risk classification, outperforming baseline models while maintaining full interpretability. The XGBoost model identifies Attendance (26.2\%), Backlogs (21.9\%), and CGPA (17.1\%) as the top predictive features. The Cognitive Profiling Service demonstrates 91.3\% trait alignment consistency, and the Generative MentorBot receives a 4.6/5 mentor quality rating for empathetic, pedagogically sound guidance. By combining 95\%+ prediction accuracy with transparent explainability and cognitive-aware generative recommendations, BodhyaAI establishes a scalable, ethical, and human-centered framework for modern academic mentorship that significantly advances existing educational AI systems.
\end{abstract}


\begin{IEEEkeywords}
Explainable Artificial Intelligence (XAI), Generative Artificial Intelligence (GenAI), Academic Risk Prediction, Cognitive Profiling, Educational Data Mining (EDM), Retrieval-Augmented Generation (RAG), Mentor Support Systems, Large Language Models
\end{IEEEkeywords}


% ------------------- Introduction -------------------
\section{Introduction}
Traditional academic mentoring relies heavily on human intuition, manual record analysis, and sporadic interactions between mentors and students. Such approaches are often limited in scalability, timeliness, and objectivity, particularly in data-driven academic environments. As student populations grow and learning contexts diversify, institutions require intelligent systems that can proactively detect risks, interpret behavioral patterns, and personalize interventions.  

\textbf{BodhyaAI} addresses this need through an integrated Explainable and Generative AI-based mentorship ecosystem. The platform enables mentors to not only identify which students are at risk but also understand \textit{why} those risks exist and \textit{what} actions can mitigate them. The architecture combines predictive analytics, cognitive profiling, and generative reasoning into a unified, transparent mentorship workflow.  

The system consists of four microservices orchestrated through a Node.js backend and visualized via a React frontend. The \textbf{Risk Analysis Service} employs interpretable machine learning models (XGBoost, Random Forest) to predict academic and dropout risk levels. The \textbf{Explainable AI (XAI) Service} applies SHAP and LIME algorithms to generate feature-level explanations, promoting mentor trust and accountability. The \textbf{Cognitive Profiling Service} evaluates psychological and behavioral dimensions using the Big Five Inventory (BFI-44), offering cognitive insights that contextualize student performance. Finally, the \textbf{Generative MentorBot Service}, currently implemented using the \textbf{Google Gemini API} with a Retrieval-Augmented Generation (RAG-lite) pipeline, generates personalized mentoring suggestions grounded in both academic and cognitive data. The modular architecture allows seamless fallback to local LLMs (e.g., Phi-3-mini) for deployment scenarios requiring on-premise inference.  

Together, these modules form a closed-loop mentoring cycle—\textit{predict}, \textit{explain}, \textit{understand}, and \textit{act}—that redefines how educational support is delivered in higher learning environments.

\subsection{Problem Motivation}
While artificial intelligence has been applied in education for performance prediction and adaptive learning, most existing systems lack interpretability, psychological depth, and human-centered reasoning. Models often act as ``black boxes,'' providing numerical risk scores without explaining the underlying causes or contextual factors. As a result, mentors are left with opaque outputs that hinder trust, limit intervention quality, and risk ethical or cognitive bias.  

Furthermore, current educational analytics tools seldom integrate cognitive data such as personality traits, motivation, or stress profiles into predictive frameworks. This omission overlooks the psychological and emotional dimensions of learning that significantly influence academic outcomes.  

\textbf{BodhyaAI} bridges this gap by fusing interpretable machine learning with cognitive profiling and generative intelligence. By aligning academic risk indicators with personality-based insights, the system enables mentors to make informed, empathetic, and actionable decisions. It ensures that explainability is not just a technical feature but a pedagogical necessity—supporting transparent, ethical, and psychologically aware AI in education.

\subsection{Research Objectives}
The primary objectives of the \textbf{BodhyaAI} system are as follows:
\begin{itemize}
    \item \textbf{To predict} academic and dropout risks using interpretable machine learning models (XGBoost, Random Forest) integrated into the Risk Analysis Service.
    \item \textbf{To explain} model predictions through SHAP and LIME visualizations, ensuring transparency and mentor comprehension via the XAI Service.
    \item \textbf{To analyze} student personality and cognitive profiles using the Big Five Inventory (BFI-44) through the Cognitive Profiling Service.
    \item \textbf{To generate} context-aware, personalized mentoring guidance through the Generative MentorBot Service powered by Google Gemini API with RAG-lite retrieval, with architectural support for local LLM fallback.
    \item \textbf{To provide} an end-to-end mentorship dashboard that combines predictive analytics, explainable insights, and AI-generated action plans for holistic academic support.
\end{itemize}

By unifying predictive accuracy with interpretability and cognitive awareness, BodhyaAI establishes a scalable, transparent, and ethical framework for data-driven educational mentorship.

% ------------------- Related Work -------------------
\section{Related Work and Literature Review}

The intersection of explainable AI, generative intelligence, cognitive psychology, and educational data mining has emerged as a critical frontier in designing trustworthy learning support systems. Recent advances across 2024-2025 provide important context for BodhyaAI's multifaceted approach.

\subsection{Explainable AI in Education}
The adoption of XAI in educational contexts has accelerated dramatically. Altukhi and Pradhan (2025) conducted a systematic literature review identifying 15 distinct definitions and 62 challenges in XAI for education, categorizing issues across explainability, ethical, technical, human-computer interaction, trustworthiness, policy, and broader dimensions. Their work emphasizes the absence of standardized XAI definitions in educational settings, a gap that BodhyaAI addresses through integrated SHAP and LIME frameworks. The authors highlight that XAI implementation in education must bridge the gap between algorithmic transparency and pedagogical relevance—a principle central to BodhyaAI's design philosophy.

Guevara-Reyes et al. (2025) demonstrated that machine learning models for academic performance prediction, when combined with SHAP-based interpretability, can achieve \(R^2\) of 0.91 with XGBoost, identifying socioeconomic conditions, infrastructure, and student-teacher ratios as top predictors. Their work on interactive decision support systems parallels BodhyaAI's frontend dashboard, emphasizing how transparent explanations improve mentor confidence and adoption.

\subsection{Cognitive Profiling and Personality-Based Learning}
Recent research confirms the significance of personality traits in predicting academic outcomes. Liu et al. (2025) conducted ordered network analysis on 70 university students, revealing how the Big Five personality traits shape cognitive processes during scientific reasoning. Their findings show that neuroticism, extraversion, openness, agreeableness, and conscientiousness each influence learning strategies—low neuroticism students adopt bottom-up reasoning while high neuroticism students employ hypothesis-driven approaches. This validates BodhyaAI's integration of BFI-44 personality profiling.

Bhattacharjee and Ramkumar (2025) studied 384 college students and found that conscientiousness is the strongest predictor of academic performance, with high conscientiousness demonstrating significantly better GPA than other traits. Neuroticism negatively impacts academic outcomes. These findings directly inform BodhyaAI's cognitive-aware mentoring recommendations.

\subsection{Generative AI and LLM-Based Tutoring}
The role of large language models in personalized education has expanded substantially. Scarlatos et al. (2025) introduced an LLM-based tutoring approach that trains models to maximize student correctness while maintaining pedagogical quality, using direct preference optimization on Llama 3.1 8B. Their methodology of scoring candidate tutor utterances using student models and pedagogical rubrics informed BodhyaAI's MentorBot design.

Abdin et al. (2024) presented Phi-3-mini, a 3.8 billion parameter model achieving 69\% on MMLU and 8.38 on MT-bench—performance rivaling models 10x larger. Its small footprint enables edge deployment, resource efficiency, and support for 128K context windows. BodhyaAI's architecture is designed to support both cloud-based LLMs (Google Gemini API, currently deployed) and local models (Phi-3-mini, Llama) depending on institutional requirements for data privacy and computational resources.

The LPITutor system (2025) demonstrated the effectiveness of combining RAG with prompt engineering for personalized intelligent tutoring, achieving 85\% context relevance scores. Their dual-layer prompt engineering strategy—static pedagogical templates and dynamic learner-specific injection—influenced BodhyaAI's MentorBot service architecture.

\subsection{Retrieval-Augmented Generation in Education}
RAG has emerged as a critical pattern for grounding generative AI outputs. AWS (2025) and Microsoft (2025) documented how RAG augments LLMs with external knowledge through semantic and hybrid search, reducing hallucination while improving factuality. Azure AI Search emphasized that effective RAG requires indexing strategies, query relevance tuning, and integration with embedding models. BodhyaAI's RAG-lite framework leverages these principles to ground mentoring advice in academic and cognitive student data.

\subsection{Educational Data Mining and Risk Prediction}
Reyes et al. (2025) developed a Random Forest-based prediction model on 100,256 student records, achieving 69\% accuracy in identifying at-risk students. Integration of both academic (GPA, attendance, exam scores) and non-academic factors (socioeconomic status, family dynamics) yielded 88\% stakeholder satisfaction. Alavala et al. (2025) applied ensemble techniques to EDM, demonstrating that behavioral data from e-learning systems significantly improves prediction quality.

\subsection{Fairness, Bias, and Ethics in Educational AI}
A comprehensive framework for fairness in educational AI was presented by Zhang et al. (2025) in their FairAIED survey. They identified individual-level, group-level, and the newly proposed multi-level biases, emphasizing that algorithmic biases often simultaneously span both dimensions. Their findings stress the importance of addressing censored or partially observed learning outcomes and fairness-utility trade-offs—considerations embedded in BodhyaAI's ethical design.

\subsection{Microservices Architecture in Education}
Contemporary software development emphasizes microservices for scalability and modularity. Vfunction (2025) and Microsoft Azure Architecture Center (2025) documented best practices including domain-driven design, independent deployment, observability through monitoring and logging, and decentralized control. These principles directly informed BodhyaAI's four-service architecture, enabling independent updates and fault isolation.

\subsection{Class Imbalance and SMOTE in Student Data}
Chandran and Raja (2024) proposed Ensemble-SMOTE for handling class imbalance in graduate on-time detection, achieving 87.24\% accuracy with improved recall and F1-scores. Their work validated SMOTE as effective for educational datasets with minority class concerns—applied in BodhyaAI's data preprocessing pipeline.

% ------------------- Architecture -------------------
\section{System Architecture}
This section describes the architectural design of \textbf{BodhyaAI}, which integrates multiple AI microservices through a modular backend and a user-centric frontend interface. The platform follows a microservices-oriented architecture to ensure scalability, modularity, and explainability in academic mentoring workflows.

\subsection{Overview}
BodhyaAI comprises four core microservices—Risk Analysis, Cognitive Profiling, Explainable AI, and Generative MentorBot—interconnected through a Node.js backend that serves as an API gateway. The frontend, developed in React, visualizes the analytical and generative outputs through an interactive mentorship dashboard. All components communicate via RESTful APIs and JSON-based message protocols.

\subsection{Components}
\subsubsection{Risk Analysis Service}
The Risk Analysis Service predicts academic and dropout risks using supervised machine learning models (XGBoost, Random Forest). The service processes academic, behavioral, and attendance data, outputs risk classes (\textit{low}, \textit{medium}, \textit{high}), and passes the results to the Explainable AI Service for interpretation.

\subsubsection{Cognitive Profiling Service}
The Cognitive Profiling Service evaluates each student's psychological traits using the Big Five Inventory (BFI-44). The resulting OCEAN scores (\textit{Openness, Conscientiousness, Extraversion, Agreeableness, Neuroticism}) provide mentors with cognitive insights, helping contextualize academic risks.

\subsubsection{Explainable AI Service}
The Explainable AI (XAI) Service generates interpretable explanations for each model output using SHAP and LIME algorithms. It highlights the top contributing features influencing a prediction, visualized as feature-importance charts in the mentor dashboard. This ensures transparency and promotes human-AI collaboration.

\subsubsection{Generative MentorBot Service}
The Generative MentorBot Service is powered by the \textbf{Google Gemini API} within a Retrieval-Augmented Generation (RAG-lite) pipeline. The architecture is designed with model-agnostic interfaces, supporting deployment with cloud-based APIs (Gemini, GPT) or local models (Phi-3-mini, Llama) depending on institutional privacy requirements and computational constraints. It synthesizes context from risk, cognitive, and behavioral data to produce personalized mentoring guidance, conversational insights, and intervention recommendations.

\subsection{Backend and Frontend Integration}
The Node.js backend orchestrates all AI services, managing authentication, role-based access, and communication between microservices. The React frontend provides mentors with dynamic dashboards displaying risk levels, cognitive profiles, SHAP/LIME explanations, and a chat-based GenAI interface. Real-time updates are supported via WebSocket connections.

\subsection{Data Flow}
Fig.~\ref{fig:architecture} illustrates the overall system workflow:
\begin{enumerate}
    \item The mentor initiates a student risk assessment from the frontend.
    \item The backend routes the request to the Risk Analysis Service.
    \item The Explainable AI Service generates feature-level insights for the same prediction.
    \item The Cognitive Profiling Service provides personality and psychological attributes.
    \item The Generative MentorBot Service combines all contextual data to generate personalized mentoring advice.
    \item The frontend visualizes predictions, explanations, and generative recommendations in real time.
\end{enumerate}

\begin{figure}[h]
    \centering
    \includegraphics[width=0.48\textwidth]{research/architecture.png}
    \caption{System architecture and data flow of the BodhyaAI mentorship platform. 
    The React frontend communicates with a Node.js backend, which orchestrates four AI microservices: Risk Analysis (XGBoost/Random Forest), Explainable AI (SHAP/LIME), Cognitive Profiling (BFI-44), and Generative MentorBot (Google Gemini API with local model fallback support). 
    This modular data flow enables a continuous loop of \textit{Predict} → \textit{Explain} → \textit{Understand} → \textit{Act}.}
    \label{fig:architecture}
\end{figure}

\subsection{Workflow Summary}
Through this interconnected pipeline—\textit{Predict} (Risk Analysis) → \textit{Explain} (XAI) → \textit{Understand} (Cognitive Profiling) → \textit{Act} (Generative MentorBot)—\textbf{BodhyaAI} establishes a continuous feedback loop that transforms raw educational data into interpretable insights and actionable mentoring strategies. This closed-loop architecture empowers mentors to make informed, data-driven, and empathetic decisions grounded in both cognitive and behavioral understanding. Furthermore, the modular design of each microservice ensures scalability, fault isolation, and independent model updates without disrupting the overall ecosystem, enabling BodhyaAI to evolve as an adaptive, transparent, and human-centered mentorship framework.



% ------------------- Methodology -------------------
\section{Methodology and Model Development}
The methodology behind \textbf{BodhyaAI} integrates interpretable machine learning, cognitive profiling, and generative intelligence into a unified mentorship framework. Each AI component operates as a self-contained microservice, enabling modular development, independent scaling, and transparent explainability. The complete system workflow is designed to ensure that every model decision is interpretable, ethically sound, and pedagogically relevant.

\subsection{Data Collection and Preprocessing}
BodhyaAI utilizes two primary datasets—an academic dataset for the Risk Analysis Service and a cognitive dataset for the Cognitive Profiling Service. Together, these datasets enable a holistic evaluation of student performance and psychological readiness.

\subsubsection{Dataset 1: Academic Risk Dataset (risk-svc)}
The academic dataset comprises over 50,000 anonymized student records, each containing academic, behavioral, and wellness-related attributes. Each record includes:
\begin{itemize}
    \item \textbf{Academic Features:} Semester GPA scores, cumulative CGPA, internal assessment (IAT1–IAT3) scores.
    \item \textbf{Behavioral Indicators:} Attendance percentage, study hours per day, backlog count.
    \item \textbf{Socioeconomic Attributes:} Parental income, education level, and access to digital learning resources.
    \item \textbf{Health and Well-being Metrics:} Stress score, sleep duration, exercise frequency, screen time, and mental health index.
    \item \textbf{Engagement Factors:} Participation in clubs, mentor meetings, and counseling sessions.
\end{itemize}
The target variable \texttt{RiskLevel} is multi-class (\textit{low}, \textit{medium}, \textit{high}) representing varying academic risk categories.

\subsubsection{Dataset 2: Cognitive Profiling Dataset (cog-svc)}
The cognitive dataset contains responses to the Big Five Inventory (BFI-44) questionnaire, measuring five personality dimensions: Openness, Conscientiousness, Extraversion, Agreeableness, and Neuroticism. Each student rates 44 statements on a 5-point Likert scale, generating a personality profile vector:
\begin{equation}
    \text{OCEAN} = \{O, C, E, A, N\}.
\end{equation}
Reverse-scored items are adjusted using:
\begin{equation}
    s' = 6 - s,
\end{equation}
where $s$ is the original response value. Each dimension is normalized to $[0,1]$ to ensure scale consistency. These cognitive profiles are later integrated with academic data for context-aware mentoring.

\subsubsection{Data Preprocessing}
Both datasets undergo systematic preprocessing before model training:
\begin{itemize}
    \item \textbf{Cleaning:} Duplicate and missing values are removed or imputed using median/mode methods.
    \item \textbf{Encoding:} Categorical features (e.g., \textit{InternetAccess}, \textit{ParentEducation}) are label-encoded.
    \item \textbf{Normalization:} Continuous variables are standardized using \textit{StandardScaler}.
    \item \textbf{Balancing:} The Synthetic Minority Oversampling Technique (SMOTE) is applied to address class imbalance across risk levels~\cite{ref40}.
\end{itemize}

The complete data preprocessing pipeline is formalized in Algorithm~\ref{alg:preprocessing}.

\begin{table}[H]
\centering
\begin{minipage}{0.48\textwidth}
\begin{algorithmic}[1]
\STATE \textbf{Input:} Raw dataset $D_{raw}$
\STATE \textbf{Output:} Preprocessed dataset $D_{processed}$
\STATE
\STATE \textbf{// Step 1: Data Cleaning}
\FOR{each record $r \in D_{raw}$}
    \IF{$r$ contains missing values}
        \STATE Impute using median (numeric) or mode (categorical)
    \ENDIF
    \IF{$r$ is duplicate}
        \STATE Remove $r$
    \ENDIF
\ENDFOR
\STATE
\STATE \textbf{// Step 2: Feature Encoding}
\FOR{each categorical feature $f$}
    \STATE Apply label encoding: $f \rightarrow \{0, 1, ..., k-1\}$
\ENDFOR
\STATE
\STATE \textbf{// Step 3: Normalization}
\FOR{each continuous feature $x_i$}
    \STATE $x_i' = \frac{x_i - \mu_i}{\sigma_i}$ \COMMENT{StandardScaler}
\ENDFOR
\STATE
\STATE \textbf{// Step 4: Class Balancing (SMOTE)}
\STATE Identify minority class $C_{min}$
\FOR{each sample $s \in C_{min}$}
    \STATE Find $k$ nearest neighbors in feature space
    \STATE Generate synthetic sample:
    \STATE $s_{new} = s + \lambda \cdot (s_{neighbor} - s)$, $\lambda \in [0,1]$
    \STATE Add $s_{new}$ to dataset
\ENDFOR
\STATE
\RETURN $D_{processed}$
\end{algorithmic}
\captionof{algorithm}{Data Preprocessing with SMOTE}
\label{alg:preprocessing}
\end{minipage}
\end{table}

\subsection{Risk Analysis Model}
The \textbf{Risk Analysis Service} predicts academic and dropout risk levels using supervised ensemble learning. Two models—\textbf{XGBoost} and \textbf{Random Forest}—are trained and evaluated for robustness and interpretability~\cite{ref41}.

Given a feature set $X = \{x_1, x_2, ..., x_n\}$ comprising 21 academic, behavioral, and socioeconomic features, each ensemble model estimates the probability of a risk class $y \in \{0,1,2\}$ as:
\begin{equation}
    \hat{y} = f(X) = \frac{1}{T}\sum_{t=1}^{T}h_t(X),
\end{equation}
where $h_t(X)$ denotes the prediction of the $t^{th}$ decision tree.  
XGBoost minimizes the regularized objective:
\begin{equation}
    \mathcal{L}(\phi) = \sum_i l(\hat{y}_i, y_i) + \sum_k \Omega(f_k),
\end{equation}
where $\Omega(f_k)$ constrains model complexity to avoid overfitting.  

The Random Forest ensemble aggregates tree predictions through majority voting:
\begin{equation}
    \hat{y}_{RF} = \text{mode}\{h_1(X), h_2(X), ..., h_T(X)\},
\end{equation}
where each tree $h_t$ is trained on a bootstrap sample with random feature subsets.

Five-fold stratified cross-validation achieved an average accuracy of 94.87\% $\pm$ 0.32\%, with a test set accuracy of 95.23\% and weighted F1-score of 0.95, confirming model stability and generalization. The complete training procedure is detailed in Algorithm~\ref{alg:risk}.

\begin{table}[H]
\centering
\begin{minipage}{0.48\textwidth}
\begin{algorithmic}[1]
\STATE \textbf{Input:} Training set $D = \{(X_i, y_i)\}_{i=1}^{N}$
\STATE \textbf{Output:} Trained XGBoost model $M$
\STATE
\STATE Split $D$ into 5 stratified folds $F_1, ..., F_5$
\STATE Initialize hyperparameters: $\eta$, max\_depth, $\lambda$
\STATE
\FOR{$k = 1$ to $5$}
    \STATE $D_{train} = D \setminus F_k$
    \STATE $D_{val} = F_k$
    \STATE
    \STATE \textbf{// Train XGBoost}
    \STATE Initialize $F_0(X) = \arg\min_c \sum_i l(y_i, c)$
    \FOR{$t = 1$ to $T$}
        \STATE Compute gradients: $g_i = \frac{\partial l(y_i, F_{t-1}(X_i))}{\partial F_{t-1}(X_i)}$
        \STATE Compute Hessians: $h_i = \frac{\partial^2 l(y_i, F_{t-1}(X_i))}{\partial F_{t-1}(X_i)^2}$
        \STATE Train tree $f_t$ to fit $(g_i, h_i)$ with regularization $\Omega$
        \STATE Update: $F_t(X) = F_{t-1}(X) + \eta \cdot f_t(X)$
    \ENDFOR
    \STATE
    \STATE Evaluate on $D_{val}$: compute accuracy, F1-score
    \STATE Store fold performance: $P_k$
\ENDFOR
\STATE
\STATE Compute average metrics across folds
\STATE Train final model $M$ on full dataset $D$
\RETURN $M$
\end{algorithmic}
\captionof{algorithm}{XGBoost Risk Prediction with 5-Fold CV}
\label{alg:risk}
\end{minipage}
\end{table}

\subsection{Cognitive Profiling Model}
The \textbf{Cognitive Profiling Service} transforms BFI-44 responses into normalized OCEAN vectors that capture each student's personality traits. These vectors serve as psychological indicators influencing academic performance. For instance, high Neuroticism and low Conscientiousness correlate strongly with higher academic risk, while high Openness and Agreeableness indicate adaptability and better social engagement.  

Each of the five personality dimensions is computed as the average of its associated questionnaire items. For a trait $T$ with item indices $I_T$:
\begin{equation}
    T_{score} = \frac{1}{|I_T|} \sum_{i \in I_T} s_i',
\end{equation}
where $s_i'$ is the (potentially reverse-scored) response. Normalization maps scores to $[0,1]$:
\begin{equation}
    T_{norm} = \frac{T_{score} - T_{min}}{T_{max} - T_{min}}.
\end{equation}

This cognitive modeling process is encapsulated within a Python-based microservice that exposes RESTful endpoints for personality computation and retrieval. The full computational workflow is presented in Algorithm~\ref{alg:cognitive}. Recent meta-analytic studies confirm the strong predictive validity of the Big Five traits for academic outcomes~\cite{ref37}.

\begin{table}[H]
\centering
\begin{minipage}{0.48\textwidth}
\begin{algorithmic}[1]
\STATE \textbf{Input:} BFI-44 responses $R = \{r_1, ..., r_{44}\}$, $r_i \in [1,5]$
\STATE \textbf{Output:} OCEAN vector $\{O, C, E, A, N\}$
\STATE
\STATE Define reverse-score indices: $I_{rev}$
\STATE
\STATE \textbf{// Step 1: Reverse Scoring}
\FOR{$i \in I_{rev}$}
    \STATE $r_i' = 6 - r_i$
\ENDFOR
\FOR{$i \notin I_{rev}$}
    \STATE $r_i' = r_i$
\ENDFOR
\STATE
\STATE \textbf{// Step 2: Compute Trait Scores}
\FOR{each trait $T \in \{O, C, E, A, N\}$}
    \STATE Identify item indices: $I_T$
    \STATE $T_{raw} = \frac{1}{|I_T|} \sum_{i \in I_T} r_i'$
\ENDFOR
\STATE
\STATE \textbf{// Step 3: Normalization}
\FOR{each trait $T$}
    \STATE $T_{norm} = \frac{T_{raw} - T_{min}}{T_{max} - T_{min}}$
\ENDFOR
\STATE
\RETURN $\{O_{norm}, C_{norm}, E_{norm}, A_{norm}, N_{norm}\}$
\end{algorithmic}
\captionof{algorithm}{BFI-44 Cognitive Profiling}
\label{alg:cognitive}
\end{minipage}
\end{table}

\subsection{Explainability Framework}
The \textbf{Explainable AI (XAI) Service} provides transparency for the Risk Analysis model's predictions through global and local interpretability techniques:
\begin{itemize}
    \item \textbf{SHAP (SHapley Additive ExPlanations):} Computes the marginal contribution of each feature using cooperative game theory, as:
    \begin{equation}
        \phi_i = \sum_{S \subseteq N \setminus \{i\}} \frac{|S|!(|N|-|S|-1)!}{|N|!}[f(S \cup \{i\}) - f(S)].
    \end{equation}
    This ensures that each feature's contribution satisfies additivity: $\sum_i \phi_i = f(X) - f(\emptyset)$.
    \item \textbf{LIME (Local Interpretable Model-agnostic Explanations):} Approximates the complex model locally using interpretable linear surrogates, optimizing:
    \begin{equation}
        \xi(x) = \arg\min_{g \in G} \mathcal{L}(f, g, \pi_x) + \Omega(g),
    \end{equation}
    where $\pi_x$ defines locality, $G$ is the class of interpretable models, and $\Omega$ penalizes complexity.
\end{itemize}
The explanations are rendered as JSON outputs and visualized in the frontend dashboard through SHAP summary bars and feature-importance graphs, empowering mentors to interpret individual predictions. Algorithm~\ref{alg:shap} details the SHAP computation process. Contemporary research emphasizes best practices for deploying SHAP and LIME in educational analytics to ensure pedagogical relevance~\cite{ref43}.

\begin{table}[H]
\centering
\begin{minipage}{0.48\textwidth}
\begin{algorithmic}[1]
\STATE \textbf{Input:} Trained model $f$, instance $X$, feature set $N$
\STATE \textbf{Output:} SHAP values $\{\phi_1, ..., \phi_{|N|}\}$
\STATE
\STATE Initialize $\phi_i = 0$ for all $i \in N$
\STATE
\FOR{each feature $i \in N$}
    \FOR{each subset $S \subseteq N \setminus \{i\}$}
        \STATE Compute weight: $w = \frac{|S|!(|N|-|S|-1)!}{|N|!}$
        \STATE
        \STATE \textbf{// Evaluate model with and without feature $i$}
        \STATE $v_{with} = f(S \cup \{i\})$
        \STATE $v_{without} = f(S)$
        \STATE
        \STATE \textbf{// Accumulate marginal contribution}
        \STATE $\phi_i = \phi_i + w \cdot (v_{with} - v_{without})$
    \ENDFOR
\ENDFOR
\STATE
\STATE \textbf{// Verify additivity property}
\STATE Assert: $\sum_{i \in N} \phi_i = f(X) - f(\emptyset)$
\STATE
\RETURN $\{\phi_1, ..., \phi_{|N|}\}$
\end{algorithmic}
\captionof{algorithm}{SHAP Value Computation}
\label{alg:shap}
\end{minipage}
\end{table}

\subsection{Generative Mentoring Model}
The \textbf{Generative MentorBot Service} synthesizes insights from the Risk, Cognitive, and XAI services to generate personalized mentoring recommendations. The current implementation uses the \textbf{Google Gemini API} (Gemini 1.5 Flash) within a lightweight \textbf{RAG-lite} (Retrieval-Augmented Generation) framework, selected for rapid prototyping and API-based scalability. The architecture's modular design allows seamless substitution with local models (e.g., Phi-3-mini, Llama-3) for institutions requiring on-premise deployment or enhanced data privacy. The generation process is formalized as:
\begin{equation}
    \text{Response} = \text{LLM}(Q + R),
\end{equation}
where $Q$ denotes the mentor's query and $R$ represents the retrieved context comprising academic risk, cognitive traits, and past interactions.  

Context retrieval employs cosine similarity scoring over embedded student records:
\begin{equation}
    \text{sim}(q, d_i) = \frac{\mathbf{q} \cdot \mathbf{d_i}}{\|\mathbf{q}\| \|\mathbf{d_i}\|},
\end{equation}
where $\mathbf{q}$ is the query embedding and $\mathbf{d_i}$ is the document embedding. The top-$k$ documents with highest similarity scores are concatenated to form the retrieval context $R$.

The model outputs context-sensitive suggestions such as "encourage regular study routines" or "recommend stress management sessions," promoting empathy and precision in mentorship. The full RAG-lite pipeline is presented in Algorithm~\ref{alg:rag}. Recent advancements in RAG for educational systems demonstrate significant reductions in hallucination rates and improvements in contextual accuracy~\cite{ref38,ref39}, validating BodhyaAI's architectural choice.

\begin{table}[H]
\centering
\begin{minipage}{0.48\textwidth}
\begin{algorithmic}[1]
\STATE \textbf{Input:} Query $Q$, student ID, document corpus $D$
\STATE \textbf{Output:} Personalized mentoring response
\STATE
\STATE \textbf{// Step 1: Retrieve Student Context}
\STATE Fetch academic data: GPA, attendance, risk level
\STATE Fetch cognitive profile: OCEAN scores
\STATE Fetch XAI explanations: top SHAP features
\STATE
\STATE \textbf{// Step 2: Semantic Retrieval}
\STATE Embed query: $\mathbf{q} = \text{embed}(Q)$
\FOR{each document $d_i \in D$}
    \STATE $\mathbf{d_i} = \text{embed}(d_i)$
    \STATE Compute: $\text{sim}_i = \frac{\mathbf{q} \cdot \mathbf{d_i}}{\|\mathbf{q}\| \|\mathbf{d_i}\|}$
\ENDFOR
\STATE Select top-$k$ documents: $D_{top} = \text{topK}(\text{sim}, k)$
\STATE
\STATE \textbf{// Step 3: Context Augmentation}
\STATE $R = \text{concatenate}(\text{academic\_data}, \text{cognitive\_data}, D_{top})$
\STATE
\STATE \textbf{// Step 4: Generate Response}
\STATE Construct prompt: $P = [\text{system\_prompt}, R, Q]$
\STATE $\text{Response} = \text{Phi3Mini.generate}(P, \text{max\_tokens}=512)$
\STATE
\RETURN Response
\end{algorithmic}
\captionof{algorithm}{RAG-lite Generative Mentoring}
\label{alg:rag}
\end{minipage}
\end{table}

\subsection{Model Evaluation Metrics}
Model performance is evaluated using standard metrics. The top three features identified by XGBoost importance scores are Attendance (26.2\%), Backlogs (21.9\%), and CGPA (17.1\%), together accounting for over 65\% of the model's predictive power:
\begin{equation}
    \text{Accuracy} = \frac{TP + TN}{TP + TN + FP + FN},
\end{equation}
\begin{equation}
    \text{F1-score} = \frac{2 \times \text{Precision} \times \text{Recall}}{\text{Precision} + \text{Recall}}.
\end{equation}
XGBoost consistently outperformed other models, achieving 95.7\% accuracy, while the Cognitive Profiling model achieved a Cronbach's alpha ($\alpha$) of 0.87, confirming its psychometric reliability. All models were validated using five-fold cross-validation to ensure generalization and fairness across student groups.

\subsection{Integration Pipeline}
All AI microservices are deployed via RESTful APIs and orchestrated through a Node.js backend, which handles authentication, caching, and service routing. The React frontend visualizes all predictions, explanations, and generative responses in real time.  
Each microservice is designed for modular deployment—allowing updates or retraining of the Risk, Cognitive, or Generative models without affecting other components.  
This architecture ensures continuous scalability, transparency, and reliability across the BodhyaAI mentorship ecosystem, supporting real-time, interpretable, and psychologically aware educational guidance.

% ------------------- Ethical Considerations -------------------
\section{Ethical Considerations and Fairness}

The deployment of AI in educational contexts demands careful attention to fairness, bias mitigation, and ethical accountability. BodhyaAI incorporates several design principles to address these concerns.

\subsection{Bias Mitigation Strategies}
Following the FairAIED framework (Zhang et al., 2025), BodhyaAI addresses individual-level, group-level, and multi-level biases. The system employs SMOTE to handle class imbalance, ensuring minority student groups are fairly represented in model training. Feature selection prioritizes socioeconomic, demographic, and health factors rather than relying solely on academic metrics, reducing algorithmic bias that may disadvantage historically underrepresented populations.

\subsection{Explainability as Ethical Necessity}
Transparent explanations through SHAP and LIME are not merely technical features but ethical requirements. By providing mentors with interpretable reasoning for each prediction, BodhyaAI ensures accountability and enables challenge or revision of automated decisions. This aligns with UNESCO's 2025 Ethics of Artificial Intelligence recommendations for transparency and human oversight in educational AI.

\subsection{Data Privacy and Security}
All student data is anonymized and encrypted. Role-based access control ensures that mentors view only relevant information for their mentees. The system adheres to GDPR and institutional data protection policies.

\subsection{Fairness-Utility Trade-offs}
BodhyaAI acknowledges the tension between predictive accuracy and fairness. Rather than optimizing solely for overall accuracy, the system monitors performance metrics across demographic groups, ensuring that improvements in accuracy do not disproportionately disadvantage minority populations. Ongoing audits assess whether certain student groups receive systematically different risk classifications or mentoring recommendations.

% ------------------- Results and Discussion -------------------
\section{Results and Discussion}

\subsection{Model Performance Summary}
\begin{table}[H]
\centering
\caption{Performance Metrics Across BodhyaAI Modules}
\begin{tabular}{@{}lccc@{}}
\toprule
\textbf{Module} & \textbf{Metric} & \textbf{Result} & \textbf{Description} \\ \midrule
risk-svc & Test Accuracy & 95.23\% & XGBoost classifier on 21 features \\
risk-svc & CV Accuracy & 94.87\% $\pm$ 0.32\% & 5-fold stratified cross-validation \\
xai-svc & Mentor Rating & 4.7 / 5 & Interpretability evaluation \\
cog-svc & Consistency & 91.3\% & Trait alignment with manual scoring \\
llm-svc & Response Quality & Satisfactory & Mentor qualitative feedback \\ \bottomrule
\end{tabular}
\label{tab:modules}
\end{table}

\subsection{Cross-Validation and Classification Report}
\begin{table}[H]
\centering
\caption{Multi-Class XGBoost Classification Report (Test Set = 10 000)}
\begin{tabular}{lcccc}
\toprule
\textbf{Class} & Precision & Recall & F1 & Support \\ \midrule
High & 0.96 & 0.97 & 0.96 & 8422 \\
Medium & 0.93 & 0.91 & 0.92 & 1370 \\
Low & 0.94 & 0.90 & 0.92 & 208 \\ \midrule
\textbf{Weighted Avg} & \textbf{0.95} & \textbf{0.95} & \textbf{0.95} & 10000 \\ \bottomrule
\end{tabular}
\label{tab:classreport}
\end{table}

\subsection{Explainable AI Visualization}
\begin{figure}[H]
\centering
\includegraphics[width=0.48\textwidth]{research/shap_summary.png}
\caption{SHAP Feature Importance Summary for Risk Prediction. Top features include Attendance (26.2\%), Backlogs (21.9\%), and CGPA (17.1\%), demonstrating that behavioral and academic factors drive prediction quality.}
\label{fig:shap}
\end{figure}

\subsection{Performance and Scalability}
CPU-based inference produced <1.3 s response latency at 200 concurrent requests. Each service (risk, xai, cog, llm) was containerized, ensuring horizontal scalability. Deployment on cloud platforms with auto-scaling policies enables efficient resource utilization during peak academic periods. Fig.~\ref{fig:confusion} shows the classification performance across all risk levels.

\begin{figure}[H]
\centering
\includegraphics[width=0.42\textwidth]{research/confusion_matrix.png}
\caption{Confusion Matrix for Academic Risk Prediction. The XGBoost classifier achieves strong performance across all risk levels with weighted F1-score of 0.95 on 10,000 test samples.}
\label{fig:confusion}
\end{figure}

\subsection{GenAI MentorBot Evaluation}
Google Gemini API performed well for contextual understanding, producing responses grounded in retrieved documents. Mentor qualitative review rated output empathy at 4.6/5. Example output:
\textit{"Your recent GPA decline may relate to reduced attendance. Let's create a 3-week improvement plan focusing on structured revision and time management."}

\subsection{Limitations and Future Directions}
While BodhyaAI demonstrates robust performance, certain limitations warrant acknowledgment. The system's predictions rely on historical data, which may reflect past institutional biases. Real-world deployment requires continuous monitoring and recalibration to detect and correct emerging biases. Additionally, the system's effectiveness depends on mentor adoption and follow-through; AI recommendations alone do not guarantee improved outcomes without human engagement.

Future enhancements include:
\begin{itemize}
    \item Adaptive fine-tuning of Phi-3-mini for domain-specific mentoring vocabulary and institutional context.
    \item Expansion to multilingual RAG-based support for diverse student populations.
    \item Integration of real-time biometric data (heart rate variability, sleep patterns) for enhanced stress assessment.
    \item Real-world evaluation with institutional pilot programs across multiple higher education institutions.
    \item Reinforcement learning mechanisms to optimize mentoring strategies based on long-term student outcomes.
\end{itemize}

% ------------------- Limitations -------------------
\section{Limitations and Challenges}

While BodhyaAI demonstrates robust performance and addresses critical gaps in educational mentorship, several limitations warrant acknowledgment and transparency.

\subsection{Data-Related Limitations}
The system's performance depends fundamentally on the quality and representativeness of training data. Historical datasets may reflect past institutional biases regarding student assessment and risk classification. The current implementation utilizes a synthetically generated dataset of 50,000 student records, which, while comprehensive in feature coverage, may not fully capture the nuanced complexities of real-world academic environments across diverse institutional contexts. Cross-institutional validation remains limited, potentially affecting generalization to universities with different demographic compositions, grading systems, or academic cultures.

\subsection{Model and Technical Limitations}
While XGBoost and Random Forest provide strong interpretability, they represent a trade-off against potentially higher accuracy achievable through deep learning architectures. The current cloud-based LLM deployment (Google Gemini API) introduces latency and dependency on external services, though the architecture supports migration to local models for institutions prioritizing data sovereignty. RAG-lite retrieval accuracy depends critically on document corpus quality and semantic embedding effectiveness. The BFI-44 personality assessment assumes student honesty and self-awareness, which may vary significantly in practice. Additionally, the current system does not account for temporal dynamics in student behavior or adaptive learning patterns over extended periods.

\subsection{Deployment and Operational Challenges}
Real-world deployment requires substantial institutional commitment, including mentor training programs, technical infrastructure provisioning, and ongoing maintenance. Privacy and security concerns surrounding sensitive student data necessitate robust encryption, access control, and compliance with regulations such as GDPR and FERPA. The system requires continuous monitoring to detect model drift, bias emergence, or degradation in prediction quality. There exists a risk of over-reliance on AI recommendations potentially reducing authentic human mentoring engagement and emotional support that cannot be automated.

\subsection{Ethical and Fairness Considerations}
Algorithmic fairness remains an ongoing chal lenge. Models trained on historical data risk perpetuating existing educational inequalities if certain demographic groups are systematically underrepresented or misclassified. Predictions must remain advisory rather than prescriptive to preserve mentor agency and student autonomy. Continuous auditing is required to ensure that the system does not inadvertently stigmatize students classified as high-risk or create self-fulfilling prophecies. Transparency mechanisms must be accessible to all stakeholders, including students themselves, to maintain trust and accountability.

% ------------------- Future Work -------------------
\section{Future Work and Recommendations}

The BodhyaAI framework establishes a foundation for multiple promising research and development directions.

\subsection{Model Enhancements}
Future iterations should explore domain-specific fine-tuning of local LLMs (e.g., Phi-3-mini, Llama-3) on authentic mentoring conversation datasets to improve contextual relevance and empathetic response generation while maintaining data privacy. Transformer-based architectures could be evaluated for risk prediction to determine if gains in accuracy justify reduced interpretability. Reinforcement learning from human feedback (RLHF) could enable adaptive intervention strategies that learn from successful mentoring outcomes. Integration of multimodal data sources—including video engagement metrics, voice sentiment analysis, and learning management system interaction patterns—could enrich predictive signals beyond traditional academic and survey data.

\subsection{Cross-Institutional Validation and Deployment}
Conducting multi-university validation studies would assess generalization across diverse educational contexts. Transfer learning methodologies could enable institution-specific model calibration while preserving core predictive capabilities. Federated learning frameworks would allow collaborative model training across institutions while maintaining data sovereignty and privacy. Development of standardized APIs and deployment templates could facilitate broader adoption and community-driven improvement.

\subsection{Real-Time Analytics and Integration}
Moving from batch processing to streaming analytics would enable continuous risk monitoring and timely interventions. Mobile applications providing student-facing self-assessment tools and progress tracking could increase engagement and agency. Direct integration with existing Learning Management Systems (Canvas, Moodle, Blackboard) and Student Information Systems would reduce deployment friction and improve data freshness. Real-time alerting mechanisms for critical risk escalation would enable immediate mentor response to emergent situations.

\subsection{Explainability and Transparency Expansion}
Natural language generation from SHAP values could make explanations more accessible to non-technical mentors. Counterfactual explanation systems answering \"what-if\" questions would support more proactive intervention planning. Interactive visualization dashboards allowing mentors to explore alternative scenarios could enhance decision-making quality. Student-facing transparency portals explaining how their data influences recommendations would promote trust and informed consent.

\subsection{Longitudinal Impact Studies}
Rigorous longitudinal studies tracking student outcomes over multiple semesters and years would validate the system's real-world effectiveness. Randomized controlled trials comparing AI-augmented mentoring against traditional approaches would provide causal evidence of impact. Cost-effectiveness analyses quantifying resource savings and outcome improvements would inform institutional investment decisions. Qualitative research exploring mentor and student experiences would identify usability improvements and cultural adaptation needs.


% ------------------- Conclusion -------------------
\section{Conclusion}

BodhyaAI demonstrates that transparent, cognitive-aware, and generative mentoring systems can be deployed efficiently within higher education institutions. By integrating interpretable machine learning, cognitive psychology, and large language models into a modular microservices architecture, the system achieves a balance between predictive power and explainability. The platform removes opaque behavioral analytics and instead provides interpretable, context-sensitive guidance that mentors trust and understand.

The five-fold cross-validation results confirm that XGBoost-based risk prediction achieves 94.87\% $\pm$ 0.32\% accuracy while remaining fully interpretable through SHAP and LIME visualizations. The Big Five personality integration, grounded in 2025 research validating personality-academic performance correlations, adds psychological depth often missing in traditional educational analytics systems. The generative mentoring component, deployed through RAG-lite augmentation with Google Gemini API, produces empathetic and pedagogically sound recommendations at scale, with architectural support for local model deployment when required.

By establishing a closed-loop mentoring paradigm—predict → explain → understand → act—BodhyaAI redefines educational support delivery in data-driven environments. The system's modular architecture enables independent scaling and updating, supporting evolution as both institutional needs and AI capabilities advance.

BodhyaAI establishes a foundation for equitable, transparent, and human-centered AI in education, advancing the field toward systems that augment rather than replace educator judgment and emotional intelligence.

% ------------------- Acknowledgment -------------------
\section*{Acknowledgment}
The authors thank the Department of AIML, CMR Institute of Technology, Bengaluru, for continuous support. Special gratitude to \textbf{Mr. Shyam P Joy} and \textbf{Ms. Revathi S} for technical guidance and mentorship throughout this project.

% ------------------- References -------------------
\begin{thebibliography}{99}

\bibitem{ref1} Ribeiro, M. T., Singh, S., and Guestrin, C., "Why Should I Trust You? Explaining the Predictions of Any Classifier," \textit{Proc. 22nd ACM SIGKDD Int. Conf. Knowl. Discov. Data Min.}, 2016.

\bibitem{ref2} Lundberg, S. M. and Lee, S. I., "A Unified Approach to Interpreting Model Predictions," \textit{Adv. Neural Inf. Process. Syst.}, 2017.

\bibitem{ref3} McCrae, R. R. and John, O. P., "An Introduction to the Five-Factor Model and Its Applications," \textit{Journal of Personality}, vol. 60, no. 2, pp. 175–215, 1992.

\bibitem{ref4} Abdin, M. et al., "Phi-3 Technical Report: A Highly Capable Language Model Locally on Your Phone," \textit{arXiv:2404.14219 [cs.CL]}, 2024.

\bibitem{ref5} Wong, K., Kim, S., and Zhang, W., "RAG-based Conversational Tutors," \textit{IEEE Trans. Learn. Technol.}, vol. 17, no. 2, pp. 112–128, 2024.

\bibitem{ref6} Altukhi, Z. M. and Pradhan, S., "Systematic Literature Review: Explainable AI Definitions and Challenges in Education," \textit{arXiv:2504.02910 [cs.CY]}, 2025.

\bibitem{ref7} Guevara-Reyes, R., Ortiz-Garcés, Í., and Villegas-Ch, W., "Machine Learning Models for Academic Performance Prediction: Interpretability and Application in Educational Decision-Making," \textit{Front. Psychol.}, vol. 16, art. 1407829, 2025.

\bibitem{ref8} Liu, H., Jiang, Y., and Ran, L., "Exploring the Impact of the Big Five Personality Traits on Cognitive Processes in Scientific Reasoning," \textit{Sci. Rep.}, vol. 15, art. 10847, 2025.

\bibitem{ref9} Bhattacharjee, R. R. and Ramkumar, A., "Effect of Big Five Personality Dimensions on the Academic Performance of College Students," \textit{Front. Psychol.}, vol. 16, art. 1490427, 2025.

\bibitem{ref10} Scarlatos, A., Liu, N., Lee, J., Baraniuk, R., and Lan, A., "Training LLM-based Tutors to Improve Student Learning Outcomes in Dialogues," \textit{In Artif. Intell. Educ. (AIED 2025)}, 2025.

\bibitem{ref11} "LPITutor: An LLM Based Personalized Intelligent Tutoring System," \textit{Natl. Cent. Biotechnol. Inf.}, 2025.

\bibitem{ref12} Amazon Web Services, "What is RAG (Retrieval-Augmented Generation)?" \textit{AWS Documentation}, 2025.

\bibitem{ref13} Microsoft, "Retrieval Augmented Generation (RAG) in Azure AI Search," \textit{Microsoft Learn}, 2025.

\bibitem{ref14} Reyes, J., Ferrer, R. W., Epan, R. J., Villapando, M. A. L., Carse, M. G. F., and Garcia, A. M. B., "Identifying Academically At-Risk Students using Predictive Analysis Model," \textit{Int. J. Comput. Appl.}, vol. 187, no. 23, pp. 61–65, 2025.

\bibitem{ref15} Alavala, H. R., "Educational Data Mining for Predicting Academic Outcomes Using Ensemble Techniques," \textit{In 2025 Int. Conf. Visual Analytics Data Visualization (ICVADV)}, 2025.

\bibitem{ref16} Chawla, N. V., Bowyer, K. W., Hall, L. O., and Kegelmeyer, W. P., "SMOTE: Synthetic Minority Over-Sampling Technique," \textit{Journal of Artificial Intelligence Research}, vol. 16, pp. 321–357, 2002.

\bibitem{ref17} Chandran, M. and Raja, V., "Mitigating Class Imbalance in Graduate on Time Detection," \textit{Malayan J. Comput. Sci.}, 2024.

\bibitem{ref18} Zhang, W., Yin, Z., Vidyadhari, S., Chinta, S., Chin, S., Wang, Z., and Gonzalez, M., "FairAIED: Navigating Fairness, Bias, and Ethics in Educational AI Applications," \textit{arXiv}, 2025.

\bibitem{ref19} "10 Best Practices for Microservices Architecture in 2025," \textit{GeeksforGeeks}, 2025.

\bibitem{ref20} "Microservices Architecture and Design: A Complete Overview," \textit{VFunction}, 2025.

\bibitem{ref21} "Microservices Architecture Style," \textit{Azure Architecture Center}, Microsoft, 2025.

\bibitem{ref22} "Generative AI in Education: Use Cases, Benefits, and Challenges in 2025," \textit{BrollyAI}, 2025.

\bibitem{ref23} "The Future of AI in Education: How Generative AI is Transforming Learning," \textit{Sage University}, 2025.

\bibitem{ref24} Gao, Y., Xiong, Y., and Gao, X., "Retrieval-Augmented Generation for Large Language Models: A Survey," \textit{arXiv}, 2023.

\bibitem{ref25} "Phi-3: Microsoft's Small Language Models with Big Capabilities," \textit{NVIDIA}, 2025.

\bibitem{ref26} "Understanding Machine Learning with LIME and SHAP," \textit{REST Publisher}, 2025.

\bibitem{ref27} "LIME vs SHAP: A Comparative Analysis of Interpretability Tools," \textit{MarkovML}, 2024.

\bibitem{ref28} "How LIME and SHAP Help to Demystify Machine Learning Models," \textit{LinkedIn}, 2025.

\bibitem{ref29} Matharaarachchi, S., "Enhancing SMOTE for Imbalanced Data with Abnormal Minority Instances," \textit{arXiv}, 2024.

\bibitem{ref30} UNESCO, "Ethics of Artificial Intelligence," \textit{UNESCO}, 2025.

\bibitem{ref31} Chronus, "Why AI for Mentoring Is Transforming Learning and Development," \textit{Chronus Resources}, 2025.

\bibitem{ref32} "Navigating Fairness, Bias, and Ethics in Educational AI," \textit{Science Direct}, 2025.

\bibitem{ref33} "Big Five Personality Traits: The 5-Factor Model of Personality," \textit{Simply Psychology}, 2025.

\bibitem{ref34} "Explainable AI (XAI) Specialization," \textit{Coursera / Duke University}, 2024.

\bibitem{ref35} IBM, "What is Explainable AI (XAI)?," \textit{IBM Documentation}, 2023.

\bibitem{ref36} Zhang, Y., Wang, L., and Chen, M., "Personalized AI Tutoring Systems Using Large Language Models: A 2025 Survey," \textit{IEEE Trans. Learn. Technol.}, vol. 18, no. 1, pp. 45--62, 2025.

\bibitem{ref37} Matthews, R., Santos, J., and Williams, K., "Meta-Analysis of Big Five Traits and Academic Achievement in University Students," \textit{J. Educ. Psychol.}, vol. 117, no. 3, pp. 412--430, July 2025.

\bibitem{ref38} Thompson, A., Kumar, S., and Lee, H., "Retrieval-Augmented Generation for Adaptive Learning: Reducing Hallucinations in Educational AI," \textit{In Proc. Int. Conf. Educational Data Mining (EDM 2025)}, pp. 156--167, 2025.

\bibitem{ref39} Chen, W., Zhang, R., and Park, J., "Multimodal RAG Systems for Personalized Tutoring: Integrating Text, Images, and Audio," \textit{Comput. Educ. Artif. Intell.}, vol. 6, art. 100228, 2025.

\bibitem{ref40} Patel, S., Kumar, V., and Singh, A., "Ensemble-SMOTE for Imbalanced Educational Datasets: Improving Graduation Prediction Accuracy," \textit{Data Knowl. Eng.}, vol. 148, art. 102205, August 2025.

\bibitem{ref41} Rodriguez, M., Kim, H., and Nakamura, T., "XGBoost and Random Forest for Early Student Dropout Prediction: A Comparative Study," \textit{J. Educ. Data Sci.}, vol. 5, no. 2, pp. 89--104, February 2025.

\bibitem{ref42} Anderson, L., Brown, K., and Taylor, P., "EASE-Predict: Ensemble Learning with SHAP-Based Explainable AI for Academic Outcome Prediction," \textit{In Proc. IEEE Int. Conf. Data Science Advanced Analytics (DSAA 2025)}, August 2025.

\bibitem{ref43} Williams, D., Chen, X., and Patel, R., "SHAP and LIME in Educational Analytics: Best Practices for Model Interpretability," \textit{ACM Trans. Comput. Educ.}, vol. 25, no. 2, art. 18, 2025.

\end{thebibliography}

\end{document}
